The method has five steps. First, we extracted the text layer of the archived Azuonye PDF and rendered the appendix chart page. Second, we compared that page with the archived chart transcription. Third, we wrote each count-bearing observation to \texttt{corpus/pagc\_inventory\_observations.jsonl} with source path, locator, abstraction level, confidence, and contradiction status. Fourth, we applied a claim gate with statuses \texttt{READY}, \texttt{REMOVE}, \texttt{REWRITE\_AS\_LIMITATION}, and \texttt{SPECULATIVE\_ONLY}. Fifth, we generated paper tables from the machine-readable corpus using \texttt{scripts/generate\_paper\_assets.py}.

This protocol treats the chart as primary for visible row and column counts, while treating phonemic splitting and matrix normalization as derived operations. It also preserves negative results: claims removed by the gate remain visible in the audit rather than silently disappearing.
